% Part: sets-functions-relations
% Chapter: infinite
% Section: dedekinds-proof
%
\documentclass[../../../include/open-logic-section]{subfiles}

\begin{document}

\olfileid[fr]{sfr}{infinite}{dedekindsproof}

\olsection[La « preuve » de Dedekind]{La « preuve » de Dedekind
de l'existence d'un ensemble infini}

Dans ce chapitre, nous avons présenté un traitement ensembliste des
entiers naturels à l'aide des algèbres de Dedekind. Dans la
\olref[sfr][arith][ref]{sec}, nous avons réfléchi, entre autres,
à la portée philosophique de l'arithmétisation de l'analyse.\footnote{Note
de l'édition française : le renvoi incomplet de la source est rétabli
vers la section de réflexions philosophiques du chapitre précédent.}
Il nous faut maintenant réfléchir à la portée de ce que nous avons
accompli ici.

Tout au long du \olref[sfr][arith][]{chap}, nous avons pris
les entiers naturels comme donnés, puis nous nous en sommes servis pour
construire explicitement les entiers relatifs, les rationnels et les réels.
Dans le présent chapitre, nous n'avons pas donné de construction explicite
des entiers naturels. Nous avons seulement montré qu'\emph{étant donné
un ensemble infini au sens de Dedekind, quel qu'il soit}, nous pouvons
définir un ensemble qui se comporte exactement comme nous souhaitons
que~$\Nat$ se comporte.

Nous ne pouvons donc manifestement pas prétendre avoir répondu à une
question métaphysique telle que : \emph{quels objets sont les entiers
naturels ?} Mais c'est une bonne chose. Après tout, dans la
\olref[sfr][arith][ref]{sec}, nous avons souligné qu'il serait erroné
de voir dans la définition de $\Real$ comme ensemble des coupures de
Dedekind une \emph{découverte}, plutôt qu'une convention commode.
L'observation cruciale est que les coupures de Dedekind possèdent les
propriétés mathématiques essentielles des nombres réels. De même ici,
l'observation cruciale est que \emph{toute} algèbre de Dedekind possède
les propriétés mathématiques essentielles des entiers naturels.
Dedekind a d'ailleurs précisé ce point en prouvant que toutes les algèbres
de Dedekind sont \emph{isomorphes}
(\citeyear[théorèmes 132--3]{Dedekind1888}). Il n'est donc pas surprenant
que de nombreux « structuralistes » contemporains voient en lui
un précurseur.

Nous avons en outre montré comment intégrer la théorie des entiers
naturels à une théorie naïve et simple des ensembles. Celle-ci reste
assez informelle, mais elle ne semble pas présupposer les entiers naturels.
Nous sommes donc peut-être en voie de réaliser l'ambitieux projet de
Dedekind lui-même, qu'il exposait ainsi :
\begin{quote}
En science, rien de ce qui est susceptible d'être démontré ne devrait
être tenu pour vrai sans démonstration. Si raisonnable que paraisse cette
exigence, je ne peux considérer qu'elle soit satisfaite, même dans les
méthodes les plus récentes pour établir les fondements de la science
la plus simple, à savoir cette partie de la logique qui traite de la
théorie des nombres. En parlant de l'arithmétique (de l'algèbre,
de l'analyse) comme d'une simple partie de la logique, j'entends affirmer
que je considère le concept de nombre comme entièrement indépendant
des notions ou des intuitions de l'espace et du temps ; je le considère
bien plutôt comme un produit immédiat des lois pures de la pensée.
\citep[préface]{Dedekind1888}
\end{quote}
L'idée audacieuse de Dedekind est donc la suivante. Nous venons de montrer
comment construire les entiers naturels à l'aide de la seule théorie
naïve des ensembles. Dans le \olref[sfr][arith][]{chap},
nous avons vu comment construire les réels à partir des entiers naturels
et d'un peu de théorie des ensembles. Peut-être « l'arithmétique
(l'algèbre, l'analyse) » s'avère-t-elle donc n'être « qu'une partie
de la logique », au sens étendu que Dedekind donne au mot « logique ».

Telle est l'idée. Mais arrêtons-nous un instant. Notre construction
d'une algèbre de Dedekind, qui nous tient lieu d'ensemble des entiers
naturels, est subordonnée à l'existence d'un ensemble infini au sens
de Dedekind. Il suffit de revenir au
\olref[sfr][infinite][dedekind]{thm:DedekindInfiniteAlgebra}
pour le constater. Si l'existence d'un ensemble infini au sens de
Dedekind ne peut être établie par la « logique » ou les « lois pures
de la pensée », le projet ne peut aboutir.

L'existence d'un ensemble infini au sens de Dedekind \emph{peut-elle}
donc être établie par les « lois pures de la pensée » ? Voici la
tentative de Dedekind :
\begin{quote}
Mon propre domaine de pensées, c'est-à-dire la totalité $S$ de toutes
les choses qui peuvent être des objets de ma pensée, est infini.
En effet, si $s$ désigne un élément de~$S$, la pensée $s'$ selon
laquelle~$s$ peut être un objet de ma pensée est elle-même un élément
de~$S$. Si nous la considérons comme une image $\phi(s)$ de
l'élément~$s$, alors \dots~$S$ est infini [au sens de Dedekind],
ce qu'il fallait démontrer.
\citep[\S66]{Dedekind1888}
\end{quote}
Il est assez stupéfiant de trouver un tel passage au milieu d'un ouvrage
composé pour une grande part de démonstrations mathématiques d'une
grande rigueur. Deux remarques s'imposent.

Premièrement, cette « preuve » n'a guère ce que nous reconnaîtrions
aujourd'hui comme un caractère « mathématique ». Elle parle d'objets
psychologiques, les pensées, et encore ne s'agit-il que de pensées
\emph{possibles}.

Deuxièmement, du moins dans la présentation que nous avons donnée des
algèbres de Dedekind, cette « preuve » présente une lacune technique
précise. Si l'argument de Dedekind réussit, il établit seulement qu'il
existe une infinité de choses, plus précisément une infinité de pensées.
Mais Dedekind doit aussi nous donner une raison de considérer~$S$ comme
un seul \emph{ensemble}, ayant une infinité d'!!{element}s, plutôt que
de penser~$S$ comme \emph{des choses}, au pluriel.

Le fait que Dedekind n'ait pas vu de lacune ici pourrait suggérer que
son emploi du mot « totalité » ne correspond pas exactement à
\emph{notre} emploi du mot « ensemble ».\footnote{Nous avons d'ailleurs
d'autres raisons de le penser ; voir \citet[p.~23]{Potter2004}.}
Cela n'aurait toutefois rien de très surprenant. Le projet que nous
avons poursuivi dans les deux derniers chapitres, à savoir une
« construction » des entiers naturels, puis, à partir de ceux-ci,
une « construction » des entiers relatifs, des réels et des rationnels,
a été entièrement mené de manière naïve. Nous nous sommes accordé tel
ou tel ensemble selon les besoins, sans poser beaucoup de principes
généraux précisant exactement quels ensembles existent, et dans quels cas.
Nous savons pourtant que \emph{certains} principes généraux sont
nécessaires : sans eux, nous retomberons sur le paradoxe de Russell.

Il est temps de dépasser notre naïveté.

\end{document}
